

%قبل از شروع مسئله ترجیحاً چیزی ننویسید. اما اگر خواستید بنویسید، در قالب لیستی (فقرات و شمارش) بنویسید.
%برای اضافه کردن نمره، بعدی از شروع‌مسئله، [x] را اضافه کنید و به جای x  نمره مورد نظر را قرار دهید. برای زیرمسئله‌ها هم به همین صورت است.
%\شروع‌مسئله[10]
\شروع‌مسئله []
%\متن‌سیاه{سوال عملی اول:}\\
%\متن‌سیاه{دراپ‌تست}\\
نکته: یک رشته پرانتزگذاری معتبر است اگر و تنها اگر ویژگی زیر را داشته باشد:
\\اگر به ازای هر اندیس از این رشته، تعداد پرانتزبازها ")" از ابتدای رشته تا به اینجا بزرگتر یا مساوی تعداد پرانتزبسته‌ها "(" از ابتدای رشته تا به اینجا باشد و در نهایت در پایان رشته تعداد پرانتزبازها و بسته‌ها برابر باشد.
\\اثبات: اگر این شرط را نداشته باشیم که واضح است پرانتزگذاری معتبر نیست چراکه در یک پرانتزگذاری معتبر همواره ابتدا یک پرانتز باز می‌شود و سپس آن پرانتز را می‌بندیم و در نهایت هم همه پرانتز ها بسته می‌شوند. حال می‌ماند اثبات اینکه اگر این ویژگی را داشته باشیم حتما پرانتزگذاری معتبر است. اثبات این مورد هم با پشته‌ای که در ادامه برای حل سوال از آن کمک می‌گیریم قابل انجام است.
\\برای این سوال از یک پشته استفاده می‌کنیم. از سمت چپ رشته شروع میکنیم و اگر در حرف \متن‌لاتین{i} ام باشیم، در صورتی که حرف \متن‌لاتین{i}ام ")" باشد \متن‌لاتین{i} را به سر پشته اضافه می‌کنیم و اگر هم "("  بود سر پشته را حذف می‌کنیم. 
\\حال برای بدست آوردن بزرگ‌ترین زیررشته‌ی معتبر، به ازای هر کدام از حروف رشته محاسبه می‌کنیم اگر انتهای زیررشته‌ی مورد نظر این حرف باشد بزرگترین بازه‌ی معتبر طولش چند است. اگر این طول از طول حداکثری که تا قبل از این بدست آورده بودیم بیشتر بود طول ماکسیمم را برابر طول جدید بدست آمده قرار می‌دهیم و تعداد زیررشته ها با طول ماکسیمم تا به این جا را برابر یک قرار می‌دهیم و اگر هم طول بدست آمده با بزرگترین طول بدست آمده پیش از این حرف برابر بود تعداد زیررشته‌های با طول ماکسیمم تا به این جای کار را هرچه بود یک واحد افزایش می‌دهیم.
\\برای بدست آوردن بزرگ‌ترین زیررشته ای که انتهایش به حرف \متن‌لاتین{i}ام ختم می‌شود به دو چیز نیاز داریم (توجه :‌ فقط برای پرانتزبسته ها این مقدار را حساب میکنیم وگرنه این مقدار صفر است): اولین مورد، آخرین پرانتز‌بسته‌ای است که تا به اینجای کار به آن رسیده‌ایم و با پرانتز‌بازی \متن‌لاتین{match} نشده (یعنی هنگامی که به آن رسیدیم پشته ما خالی بوده است) که اگر حرف \متن‌لاتین{i}ام که اکنون روی آن هستیم پرانتر بسته بود و پشته‌ی ما خالی بود این مقدار را برابر همین حرف \متن‌لاتین{i}ام قرار می‌دهیم و به حرف بعدی می‌رویم چرا که این پرانتز با پرانتزی \متن‌لاتین{match} نمی‌شود و زیررشته‌ی معتبری شامل این حرف نداریم(برای اثبات از شرط گفته شده برای معتبر بودن پرانتزگذاری استفاده می‌کنیم).
اما اگر پشته خالی نبود همان طور که گفتیم ابتدا سر آن را حذف می‌کنیم. سپس اگر پشته خالی نشد، ابتدای بزرگترین زیررشته‌ی معتبری که به حرف \متن‌لاتین{i}ام ختم می‌‌شود، حرف بعدیِ سرِ پشته‌ی کنونی است. چرا که در این بازه شرط گفته شده برای معتبر بودن رشته برقرار است. پس اگر حرف \متن‌لاتین{i}ام باشیم و عدد سر پشته (بعد از حذف سر قبلی) \متن‌لاتین{j} باشد، زیررشته‌ی
\\ \متن‌لاتین{s[j + 1], s[j + 2] ..., s[i]} بزرگترین زیررشته معتبری است که انتهایش \متن‌لاتین{i} است.
\\اما اگر پس از حذف سر پشته، پشته خالی شد، شروع بزرگترین بازه‌ی معتبری که به حرف \متن‌لاتین{i}ام ختم می‌شود، اندیس \متن‌لاتین{p + 1} است به طوری که \متن‌لاتین{p} اندیس آخرین جایی است که به پرانتز بسته‌ای رسیدیم که با چیزی \متن‌لاتین{match} نشد و پشته خالی بود. زیرا در این زیررشته شرط معتبر بودن که گفته شد برقرار است. شروع این بازه قبل از اندیس \متن‌لاتین{p + 1} هم نمی‌تواند باشد زیرا در غیر اینصورت باز هم با چک کردن شرط گفته شده می‌فهمیم شرط معتبر بودن برقرار نیست.
‌\\کد نمونه (با تشکر از حمیدرضا کلباسی) :
%}
% صورت اصلی مسئله را در این‌جا بنویسید.
%اگر مسئله چند قسمت داشت به این صورت عمل کنید:
%\شروع{زیرمسائل}
%\زیرمسئله[0] % مسئلۀ ۱آ
%\زیرمسئله % مسئلۀ ۱ب
%\زیرمسئله % مسئلۀ ۱پ
%\زیرمسئله % مسئلۀ ۱ت
%\پایان{زیرمسائل}



%%%%%%%%%%%%%%%%%%%%%%%%%%%%%%%%%%%%%%%%%%%%%%%%%%%%%
% برخی دستورات مفید برای تمرین‌ها. %
%%%%%%%%%%%%%%%%%%%%%%%%%%%%%%%%%%%%%%%%%%%%%%%%%%%%%

% فرمول‌های نمونه:
%$f(x) = \Theta(x)$
%$T(x, y) \leq \log(x) + 2^y + \binom{2n}{n}$
% {\tt code\_function}


% لیست موارد بدون شماره‌گذاری:
%\شروع{فقرات}
%    \فقره مورد اول 
%        \شروع{فقرات}
%            \فقره مورد اول داخلی 
%                \شروع{فقرات}
%                    \فقره اولین مورد خیلی داخلی 
%                    \فقره دومین مورد خیلی داخلی 
%                \پایان{فقرات}
%            \فقره مورد دوم داخلی 
%        \پایان{فقرات}
%    \فقره مورد دوم 
%\پایان{فقرات}

% لیست موارد شماره‌گذاری شده:
%\شروع{شمارش}
%    \فقره یک مورد 
%    \فقره یک مورد دیگر 
%    \فقره باز هم یک مورد دیگر
%\پایان{شمارش}

% برای نوشتن یک فرمول چند خطی که یک شماره بخوره. برای حالت‌های دیگر به سایت زیر مراجعه کنید:
% https://www.overleaf.com/learn/latex/Aligning_equations_with_amsmath
%\begin{align} 
%    \begin{split}
%        (x+y)^3 &= (x+y)^2(x+y) \\
%                &= (x^2+2xy+y^2)(x+y) \\
%                &= (x^3+2x^2y+xy^2) + (x^2y+2xy^2+y^3) \\
%                &= x^3+3x^2y+3xy^2+y^3
%    \end{split}                                 
%\end{align}

% نوشتن ماتریس و آرایه دوبعدی
%\begin{align}
%    A = 
%    \begin{bmatrix}
%        A_{11} & A_{21} \\
%        A_{21} & A_{22}
%    \end{bmatrix}
%\end{align}

% اضافه کردن تصویر با شرح. برای اضافه کردن تصویر بدون شرح، فقط دستور درج‌تصویر را استفاده کنید.
%\شروع{شکل}[h!]
%\تنظیم‌ازوسط
%  \درج‌تصویر[پهنا=0.5\پهنای‌متن]{img.jpg}
%\شرح{شرح تصویر}
%\برچسب{تصویر:۱}
%\پایان{شکل}
%\پایان{زیرمسائل}

%کد زیر را ببینید.
%%\lstinputlisting{source_filename.py}
\begin{latin}
\begin{lstlisting}
input()
a = input()

stack = []
i = 0
cnt = 0
mx = 0
pv = 1
for x in a:
  i += 1
  if (x == '('):
    stack.append(pv)
    pv = i + 1
  else:
    if len(stack) == 0:
      pv = i + 1
      continue
    j = stack.pop()
    pv = j
    l = i - j
    if l == mx:
      cnt += 1
    if l > mx:
      cnt = 1
      mx = l

if mx == 0:
  print(0, 1)
else:
  print(mx + 1, cnt)
\end{lstlisting}
\end{latin}

